\chapter{Conclusion and Future Work}
\label{ch:conclusions}

\section{Introduction}
\label{sec:conclusion-introduction}

This chapter concludes the project. It summarizes the implemented system, relates the main findings to the available evidence, states the contributions and limitations, and presents practical future work. The conclusions are limited to the prototype and evaluation described in Chapters~\ref{ch:architecture}--\ref{ch:results}.

\section{Summary of the Implemented System}
\label{sec:implemented-system-summary}

A hybrid local-governance reporting prototype is implemented. Its blockchain layer is a private Geth Clique Proof-of-Authority network with chain ID 1337, three validator nodes, and one non-mining RPC node. Four smart contracts are deployed: \texttt{Reporting}, \texttt{EmergencyReporting}, \texttt{OpinionPolling}, and \texttt{AuthorityMultiSig}. Standard reports, emergency reports, community decisions, authority actions, opinion polls, and governed authority access are supported.

Citizen writes are processed by a relayer. AI moderation, IPFS upload, and blockchain submission are arranged through a Redis-backed BullMQ pipeline, and progress is sent to the client through Server-Sent Events (SSE). Accepted complaint data and images are stored in one IPFS JSON envelope, while its CID and related hashes are recorded on-chain. A Next.js dApp is provided for citizens and authorities, and a separate Vite portal is provided for the government identity simulator.

The AI layer is a centralized off-chain aggregator with safety, spam, and civic-relevance services. It is not a decentralized oracle network. CLIP-based civic image relevance is disabled. Citizen activity is represented by pseudonyms, and government-signed one-use tickets are used as nullifiers. However, the identity service is a simulator and does not generate or verify a true zero-knowledge proof.

\section{Key Findings Against the Evidence}
\label{sec:conclusion-findings}

The prototype goal of implementing a permissioned blockchain and contract-controlled civic workflow has been achieved. About 75 contract tests are present in the current inventory. The strongest recorded execution evidence has been provided for the \texttt{Reporting} and \texttt{EmergencyReporting} contracts. Selected replay, duplicate-action, access-control, relayer, AI, IPFS, and integration checks have also been recorded. One complete report workflow has been executed from simulated identity issue to report closure. These results support the conclusion that the main components can be integrated and that the tested workflow rules can be enforced.

The interface goal has been partly achieved. Citizen and authority interfaces are implemented, gas handling is hidden by the relayer, and progress is shown through SSE. However, no automated frontend tests or broad user study has been completed. Usability, accessibility, and browser compatibility have therefore not been established.

The privacy and AI goals have also been partly achieved. Direct government identity data is not stored in the contracts, one-use tickets control tested duplicate actions, and stable pseudonyms are used on-chain. Selected moderation cases have been processed before IPFS and blockchain submission. These results do not establish anonymous credentials, zero-knowledge privacy, broad AI accuracy, or decentralized oracle operation.

Recorded timing and throughput values suggest that the system is usable as a controlled prototype. However, raw benchmark artifacts, repetition counts, and several test details have not been committed. The values are therefore prototype findings with limited reproducibility and cannot be generalized to municipal scale.

\section{Project Contributions}
\label{sec:project-contributions}

The main contribution is the integration of contract-enforced civic reporting with off-chain processing in one working prototype. The project provides separate workflows for standard and emergency reports, community validation and post-resolution review, multisignature authority governance, and opinion polling. It also shows how a relayer, BullMQ, SSE, AI checks, and an IPFS complaint envelope can be combined with a small permissioned blockchain.

A second contribution is a clear separation between public records and trusted off-chain services. Report states and important actions are recorded on-chain, while identity checking, moderation, queue processing, and media storage remain off-chain. This separation makes the trust boundaries visible. A further contribution is the use of government-signed one-use tickets and pseudonyms as a practical prototype for controlled participation, without presenting them as a complete anonymity solution.

\section{Limitations}
\label{sec:conclusion-limitations}

The project is not production-ready. The identity component is a simulator, not a true anonymous-credential or zero-knowledge system. The relayer, ticket issuer, Redis instance, centralized AI aggregator, IPFS service, and validator operators remain trusted components. Stable pseudonyms and service metadata may permit linking of activity.

The AI evidence is based on small selected samples, and broad multilingual evaluation has not been completed. CLIP civic image relevance is disabled, and no human appeal process is implemented. No external smart-contract audit, formal verification, or broad penetration test has been completed.

Frontend automated tests, accessibility evaluation, a broad user study, and a production pilot are absent. Long-term IPFS persistence, Redis high availability, validator failure, and validator governance have not been established. Raw AI and performance benchmark artifacts have not been committed, so important recorded results cannot be independently reproduced from the repository alone.

\section{Prioritized Future Work}
\label{sec:future-work}

Future work should be completed in the following order.

\begin{enumerate}
    \item \textbf{Introduce real privacy credentials.} The identity simulator should be replaced by a reviewed anonymous-credential or zero-knowledge protocol. Eligibility and one-use actions should be proved without exposing a government identifier or a stable citizen key.

    \item \textbf{Secure services and keys.} Authentication and authorization should be added to internal services, especially the IPFS endpoint. Secrets should be stored in a managed key service, key rotation and recovery should be defined, and metadata retention should be minimized.

    \item \textbf{Obtain independent assurance.} The smart contracts and relayer should receive an external security audit. Important contract invariants, access rules, nullifier use, and state transitions should also be checked through formal verification where practical.

    \item \textbf{Improve AI evidence and oversight.} Representative Sinhala, Tamil, and English civic data should be collected with suitable consent and governance. Accuracy, false rejection, fairness, and robustness should be evaluated, and rejected users should be given a human review and appeal path. Image relevance should be enabled only after a separate validated evaluation.

    \item \textbf{Evaluate decentralized moderation governance.} A decentrally governed oracle option may be tested if independent operators and clear decision rules are available. It should be compared with the centralized aggregator for cost, delay, accountability, and failure handling before adoption.

    \item \textbf{Strengthen operational resilience.} Redis replication and automatic failover should be introduced. IPFS pinning, backup, restoration, and retention policies should be tested. Validator admission, removal, rotation, outage response, and recovery should be governed and documented.

    \item \textbf{Test the frontend and accessibility.} Automated unit, integration, and end-to-end tests should be added for citizen, authority, and governance paths. Mobile use, browser compatibility, keyboard access, screen-reader support, and relevant accessibility guidance should be evaluated.

    \item \textbf{Run a limited municipal pilot.} A small supervised pilot should be conducted only after the earlier security and reliability work. It should include realistic tasks, incident procedures, a System Usability Scale (SUS) study, task completion measures, interviews, and measured operational load. Findings should be used before any wider deployment is considered.
\end{enumerate}

\section{Final Conclusion}
\label{sec:final-conclusion}

The project has produced a working prototype for blockchain-supported local-governance reporting. The private PoA chain, four contracts, relayer pipeline, IPFS storage, web interfaces, identity simulator, and centralized AI moderation are integrated. Recorded tests have supported selected contract rules, security controls, service integrations, and one complete workflow.

The strongest project goals have been achieved at prototype level for blockchain deployment, contract workflows, and component integration. Interface validation, privacy, AI generalization, decentralization, security assurance, and real-world adoption have been achieved only partly or remain future work. The project therefore demonstrates technical feasibility in a controlled setting, but it does not demonstrate production readiness. Its main value is a testable foundation from which stronger privacy, security, evidence, resilience, and municipal evaluation can be developed.



% \chapter{Conclusions and Future Work}
% \label{ch:conclusions}

% This chapter brings the research, architectural design, and engineering implementation of the blockchain-based, community-assisted, privacy-preserving framework for local governance to a close. It summarizes the main research contributions and system innovations, assesses the empirical insights and technical trade-offs that were established during the development of the prototype, examines the limitations of the present implementation, and sets out the strategic paths for future academic research, extensions to the protocol, and actual deployment in real municipal settings.

% \section{Summary of Research and Engineering Contributions}
% \label{sec:summary-contributions}

% Standard centralised civic reporting systems—for example, municipal hotlines, large-scale e-governance web portals, and commercial services such as \textit{FixMyStreet} or \textit{SeeClickFix}—have systemic drawbacks. These drawbacks are opaque resolution procedures, the existence of centralised data silos, susceptibility to arbitrary administrative rejection, the absence of inter-agency accountability, and citizens' reluctance to report sensitive public infrastructure or governance failures because of fear of retaliation. This thesis presents, designs, and implements a decentralised, privacy-first civic engagement ecosystem which overcomes these problems by means of a six-layer hybrid on-chain/off-chain architecture.

% The primary research and engineering contributions of this work are summarized as follows:
% \begin{enumerate}
%     \item \textbf{Decentralized Lifecycle Accounting via Deterministic Finite State Machine (FSM):} The project developed an eight-state deterministic finite state machine (\ac{FSM}) as part of the \texttt{Reporting.sol} smart contract (states $S_0$ through $S_7$ being \texttt{PendingValidation} to \texttt{Reopened} respectively). It did so by adhering to strict Checks-Effects-Interactions (CEI) design patterns and by keeping track of a single-use cryptographic ticket. The ledger ensures that there is an immutable and tamper-evident audit trail of all civic reports by using nullifiers in persistent mappings (specifically those named \texttt{usedSubmission\allowbreak Nullifiers}, \texttt{usedValidation\allowbreak VoteNullifiers}, \texttt{usedVerification\allowbreak VoteNullifiers}, and \texttt{usedRejection\allowbreak ReviewVote\allowbreak Nullifiers}), while also providing a mathematical prevention of Sybil attacks and double-voting.
%     \item \textbf{Two-Tier Multi-Signature Governance and On-Chain Profile Directory:} To avoid giving any single authority unilateral administrative control and to promote civic transparency, access control was decentralised by means of the \texttt{AuthorityMultiSig.sol} contract, to which the ownership of \texttt{Reporting.sol} is passed when the contract is deployed. The contract sets up a two-tier system of governance: Super Admins (who are high-ranking municipal officials and vote with a dynamic majority quorum $\lfloor \texttt{superAdminCount} / 2 \rfloor + 1$) and Authority Workers (who are municipal staff responsible for carrying out physical infrastructure repairs). Each official is given an immutable on-chain identity profile (containing their \texttt{name}, \texttt{position} and \texttt{department}), so that citizens can verify exactly which institutional authority is responsible for dealing with a reported problem.
%     \item \textbf{Zero-Knowledge Citizen Authentication and Identity Decoupling:} To remove the fear among citizens of suffering retaliation without at the same time compromising the integrity of the platform, the system made use of a Zero-Knowledge Proof (\ac{ZKP}) GovID Simulator. In order to obtain a batch of randomly generated, one-off ticket identifiers ($\texttt{zkpTicketId}$), citizens carry out their authentication procedure against the official municipal registry, these identifiers being signed with the ECDSA private key of the government authority. The platform is able to achieve a full decoupling between a citizen's real-world civic identity and their on-chain reporting activities by determining deterministic citizen pseudonyms ($\text{keccak256}(\text{citizenPubKey} \parallel \text{DOMAIN\_SALT})$) on the client side.
%     \item \textbf{Gasless Relayer Middleware and Automated Orchestration:} To eliminate the obstacles associated with getting started with cryptocurrency—such as installing a wallet, obtaining gas tokens, and configuring an RPC—the architecture introduced a NestJS backend relayer. This relayer includes a Dual-Layer Cryptographic Verification Engine which checks both the ECDSA signatures of government tickets and the signatures of the citizen payloads before sending gas-sponsored transactions to the blockchain. Additionally, the relayer has two automated backend orchestration services:
%     \begin{itemize}
%         \item \textit{Authority Funding Service (\texttt{AuthorityFundingService}):} An event-driven listener which keeps a watch on \texttt{ProposalExecuted} events occurring on \texttt{AuthorityMultiSig.sol} and automatically transfers any native ETH top-ups to the Authority Worker accounts that have been newly approved, thus smoothly linking on-chain access control with off-chain operational readiness.
%         \item \textit{Background Cron Orchestration (\texttt{CronService}):} Responsible for carrying out automated tasks that keep an eye on the balances of authority wallets and carry out the daily midnight batch finalisations (\texttt{batchFinalizeVotingWindows}) for the reports relating to the expired voting phases ($S_0$, $S_4$, and $S_5$), thus avoiding workflow stagnation.
%     \end{itemize}
%     \item \textbf{Off-Chain P2P Media Storage and Multi-Oracle AI Content Moderation:} To avoid the blockchain becoming overloaded and to deal with spam, abusive language, or non-civic noise before the data is recorded on the blockchain, the design separated state management from data storage. The high-resolution photographic evidence is stored via peer-to-peer \ac{IPFS} storage (\texttt{ipfsCid}), and the text of the reports together with the images are assessed by an interconnected \ac{AI} Content Moderation Oracle Ensemble (this consisting of a Safety Oracle which uses fine-tuned \texttt{unitary/toxic-bert}, a Spam and Abuse Oracle which uses \texttt{sms-spam-detection}, a Civic Relevance Oracle which uses \texttt{sentence-transformers/all-MiniLM-L6-v2}, and an Oracle Aggregator).
%     \item \textbf{Democratic Lifecycle Oversight and Anonymous Public Consultation:} The framework gives citizens the opportunity to take part at important decision points: voting for peer validation ($S_0$) helps to filter out spam, voting for post-repair verification ($S_5$) stops authorities from closing reports that have not been resolved, and voting for rejection review ($S_4$) enables citizens to challenge cases where issues have been improperly dismissed. Moreover, the \texttt{OpinionPolling.sol} contract built into the system allows municipal authorities to carry out anonymous public opinion consultations that are resistant to Sybil attacks regarding municipal proposals and budget allocations.
%     \item \textbf{Asynchronous BullMQ Flow Queue and Real-Time SSE Notification Gateway:} To address the latency barrier inherent in multi-step blockchain submission pipelines, the backend relayer is extended with a Redis-backed BullMQ job queue. Report submissions are processed as Flow Jobs — a parent job containing three sequentially dependent child steps (AI moderation, IPFS upload, and blockchain submission). Each step uses exponential-backoff retry logic. A \ac{SSE} notification gateway streams live progress events from queue workers to the citizen's browser, giving citizens step-by-step confirmation of their submission status without requiring page refreshes or long synchronous waits. A dedicated vote queue handles community vote submissions asynchronously in the same way.
% \end{enumerate}

% \section{Synthesis of Key Findings and Engineering Validation}
% \label{sec:synthesis-findings}

% The development and evaluation of the prototype yielded several critical insights into the design of decentralized civic governance systems:

% \begin{enumerate}
%     \item \textbf{Reconciling Public Transparency with Citizen Privacy:} Establishing a balance between public transparency and citizen privacy presents a fundamental challenge in the case of public blockchain applications, since the immutable records on the chain are in conflict with requirements relating to user privacy and data protection. This study showed that by using one-off cryptographic ticket nullifiers together with ECDSA signatures issued by the government, it is possible to overcome this problem. This enables citizens to publicly demonstrate that they are entitled to report problems and vote on local issues without having to associate their transactions with a fixed wallet address or disclose any personally identifiable information (\ac{PII}).
%     \item \textbf{Middleware Abstraction as a Prerequisite for Civic Adoption:} The adoption of middleware abstraction is a necessary precondition: empirical evaluation has shown that asking citizens to interact directly with decentralised ledgers results in an insurmountable usability barrier when it comes to municipal applications. The system attains the kind of latency and usability typical of standard web applications (through the use of the Next.js presentation layer) while at the same time maintaining the cryptographic verification and censorship-resistance assurances of the underlying blockchain by delegating responsibilities such as transaction signing, gas sponsorship, and RPC connection management to an off-chain backend relayer.
%     \item \textbf{Automated Cron Finalizations Complement Lazy Evaluation:} The use of automated daily finalisations through \texttt{CronService.\allowbreak batchFinalizeVotingWindows} complements lazy evaluation: in decentralized smart contracts state transitions generally depend on 'lazy evaluation', which means that the state updates are carried out when a later user transaction interacts with the contract; but in the case of civic reporting, reports that have no voting activity by the time the phase deadline has passed may remain indefinitely in a pending state. Tests have shown that it is essential to combine lazy evaluation with daily automated batch finalisations via \texttt{CronService.\allowbreak batchFinalizeVotingWindows} if deterministic progression of the lifecycle is to be guaranteed without needing constant manual administrative intervention.
%     \item \textbf{Efficacy of Multi-Oracle Pre-Submission Filtering:} The effectiveness of using multiple oracles for pre-submission filtering was shown by the benchmarking of the \ac{AI} moderation oracle ensemble in that off-chain machine learning verification is essential for guarding private permissioned ledgers against spam, flood attacks, and irrelevant content. When toxicity and non-civic reports are filtered at the relayer gateway, this stops the ledgers from being polluted in an immutable way and makes it possible for municipal engineering teams to concentrate solely on carrying out actionable repairs to civic infrastructure.
%     \item \textbf{Asynchronous Queue Processing as a Prerequisite for Civic Web3 UX:} A key finding from the implementation is that synchronous report submission — where the citizen's HTTP request blocks until AI moderation, IPFS upload, and blockchain confirmation all complete — is not acceptable for a civic application. The total pipeline duration ranges from 10 to 25 seconds under normal conditions. The BullMQ Flow Queue integration resolves this by returning the \texttt{jobId} immediately upon request acceptance and streaming live progress via \ac{SSE}. This pattern makes the blockchain-backed system feel as responsive as a conventional web application. It also improves resilience: if IPFS or the blockchain RPC endpoint experiences a transient failure, the retry logic in the queue recovers automatically without requiring the citizen to resubmit.
% \end{enumerate}

% \section{Limitations of the Current Prototype}
% \label{sec:limitations}

% While the prototype successfully demonstrates the feasibility of a privacy-preserving, blockchain-based local governance platform, several technical and architectural limitations remain:

% \begin{enumerate}
%     \item \textbf{Simulated Zero-Knowledge Identity Infrastructure:} In the present implementation, citizen authentication and ticket signing are carried out using the \ac{ZKP} GovID Simulator, which is a \texttt{Node.js}/\texttt{Express} service that mimics a national identity registry. Even though the cryptographic ticket verification protocol accurately reproduces zero-knowledge attestation, for real-world use it is necessary to integrate with a genuine zero-knowledge proving circuit so that citizens can generate their zero-knowledge proofs locally on their own devices without having to trust a mock server.
%     \item \textbf{Co-Located Oracle Aggregation Topology:} Although the \ac{AI} Moderation Ensemble logically divides off the verification of Safety, Spam, and Civic Relevance into separate worker modules, in the present prototype these services are placed within the same administrative environment. This means that the oracle layer now functions as a verification gateway that is co-located administratively rather than as a genuinely decentralized, multi-organizational oracle network.
%     \item \textbf{Unimodal Textual Moderation Pipeline:} The current \ac{AI} content moderation pipeline makes use of unimodal natural language processing (\ac{NLP}) models when assessing textual descriptions. Although visual material is checked for adult or graphic content, the prototype does not include a multimodal reasoning model which could cross-verify whether the uploaded image matches semantically the issue category chosen by the citizen and the textual description.
%     \item \textbf{Permissioned Consortium Scale and Validator Latency:} The blockchain in question functions as an Ethereum-compatible Clique Proof-of-Authority (\ac{PoA}) network within a consortium consisting of three validators (those of the Municipal Authority, the Regional Oversight, and the Civic NGO). Although this kind of setup offers quick block finality and entails no energy cost when compared to Proof-of-Work (\ac{PoW}), it still needs to undergo more empirical stress testing in order to ensure that its network throughput and consensus stability can handle large volumes of municipal transactions (for example, thousands of citizen votes taking place at the same time during a city-wide emergency).
% \end{enumerate}

% \section{Directions for Future Research and System Extensions}
% \label{sec:future-work}

% To address the limitations of the current prototype and transition the framework toward production readiness, future research and development should focus on four key strategic directions:

% \subsection{Production Zero-Knowledge Proof (ZKP) and Self-Sovereign Identity (SSI) Protocol Integration}
% \label{subsec:future-zkp-ssi}

% Later versions should substitute the ZKP GovID Simulator with a production-grade \textbf{Self-Sovereign Identity (\ac{ssi})} framework that complies with the W3C Decentralized Identifier (DID) and Verifiable Credential (VC) standards; specifically, incorporating open-source zero-knowledge identity protocols such as Semaphore or \textbf{Semaphore} or \textbf{Iden3 / Polygon ID} would enable municipal governments to give out cryptographic residency credentials to their citizens.

% With this model, citizens would create Groth16 or PLONK zero-knowledge proofs directly in their web browser, demonstrating that they have a valid municipal credential and that they have not gone over their voting or reporting frequency limits—without having to reveal their National Identity Number, their real name, or their persistent public key to the Backend Relayer or smart contract.


% \subsection{Multimodal Vision-Language Models (VLMs) for Cross-Modal Verification}
% \label{subsec:future-vlms}

% In order to improve the accuracy of automated moderation and to identify sophisticated cases of civic spam, the AI Moderation Oracle Ensemble should be updated to include open-weight \textbf{Multimodal Vision-Language Models (VLMs)} (for example, fine-tuned LLaVA, Llama-3-Vision, or CLIP-based cross-modal attention architectures). 

% A VLM oracle worker would take in at the same time the citizen's textual description, the selected infrastructure category, and the uploaded photographic evidence (\texttt{ipfsCid}). The model would then calculate a semantic cross-modal consistency score to check whether the image shows the civic defect that was reported (for instance, verifying that a report classified as 'Road Hazard' and involving text that mentions a pothole actually includes a genuine image of a pothole, not stock internet images or a photograph that has nothing to do with the report). This cross-modal verification would greatly increase the platform's resistance to automated Sybil attacks and spam.

% \subsection{Multi-Organizational Decentralized Oracle Network (DON) Deployment}
% \label{subsec:future-don}

% In order to avoid depending on a single administrative body for off-chain content moderation, future work ought to transform the AI Oracle layer into a genuine \textbf{Decentralized Oracle Network (DON)}. Rather than running the oracle workers on a single backend infrastructure, each oracle node should be hosted by separate and independent administrative organisations—for example, local university computer science departments, independent civic \ac{NGO}s, and municipal IT departments. 

% The independent oracle operators should use a Byzantine Fault Tolerant (BFT) consensus protocol or a threshold signature scheme (such as BLS signatures), so that a civic report will only be approved for on-chain anchoring after a cryptographic supermajority (for example, $\ge 2 / 3 + 1$) of the independent nodes carries out their own machine learning inference and signs an \texttt{ACCEPT} decision. This improvement would ensure that no single government organisation or technical administrator could unilaterally censor or change civic submissions.

% \subsection{Real-World Pilot Field Study and Usability Evaluation}
% \label{subsec:future-pilot}

% The system should first be tested through a well-organized, multi-stage pilot study in a local municipal council or at a university campus in Sri Lanka. This field study must assess both the quantitative system metrics and the qualitative user experiences of three groups of stakeholders: citizens, municipal engineers (known as Authority Workers), and municipal administrators (called Super Admins). 

% The main empirical evaluations should be:
% \begin{itemize}
%     \item \textbf{Usability and Accessibility Metrics:} Measuring, using mobile and desktop browsers, the System Usability Scale (SUS) scores, task completion rates, and average time-to-report among diverse demographic groups.
%     \item \textbf{Operational Efficiency Benchmarks:} Benchmarks relating to operational efficiency: assessing the end-to-end resolution time for municipal cases—that is, comparing the time that elapses between the submission of the initial report and its verified closure when using the existing municipal reporting channels.
%     \item \textbf{Relayer Sustainability and Gas Accounting:} The assessment of relayer sustainability and gas accounting involves measuring the amount of gas used by the relayer's treasury and the rates at which the \texttt{AuthorityFundingService} is replenished under actual municipal workloads in order to achieve long-term economic sustainability.
% \end{itemize}

% \subsection{Redis Queue Resilience and High Availability}
% \label{subsec:future-redis}

% The current BullMQ queue depends on a single Redis instance. In a production deployment, this is a single point of failure for the asynchronous submission pipeline. Future work should replace the single-instance Redis configuration with a Redis Sentinel or Redis Cluster setup. Redis Sentinel provides automatic failover by monitoring a primary node and promoting a replica if the primary becomes unavailable. Redis Cluster provides horizontal partitioning and high availability across multiple nodes. Either configuration would ensure that queued jobs are not lost during infrastructure failures and that the SSE notification gateway remains operational continuously.

% \section{Final Concluding Remarks}
% \label{sec:final-remarks}

% The thesis shows that the structural weaknesses of centralized local governance reporting systems can be effectively dealt with by combining private permissioned blockchain networks, zero-knowledge privacy protocols, automated relayer orchestration, and artificial intelligence for content moderation. Through the introduction of an eight-state deterministic workflow, two-tier multi-signature access control, and Sybil-resistant community voting, the proposed framework changes civic issue reporting from an opaque and administrative process into a transparent, democratic, and mutually accountable civic partnership. 

% Even though the present prototype is at an intermediate stage of development, its architectural strength and the fact that it has been successfully verified prove that privacy-preserving, community-assisted blockchain platforms have great potential for the future of local governance and participatory civil society.